\begin{explanationbox}
\textbf{一、最直观的画面：固定斜率，把直线向上或向下平移}\par
先固定 $k$，直线 $y=kx+m$ 的倾斜方向就固定了；改变 $m$，整条直线只会\textbf{平行移动}。
\begin{itemize}
    \item 当 $m=0$ 时，直线经过椭圆中心，必定穿过椭圆，产生两个交点。
    \item 从 $m=0$ 开始增大 $m$，直线向上平移：先有两个交点；移动到\textbf{刚好碰到椭圆}的位置时只剩一个交点，即相切；再往上就没有交点。
    \item 向下平移也完全一样。由于椭圆关于原点中心对称，临界截距必定是成对的 $m=\pm\sqrt{a^2k^2+b^2}$。
\end{itemize}
因此，\textbf{同一个斜率对应两条平行切线}。判断位置关系，本质上就是比较 $|m|$ 与临界值 $\sqrt{a^2k^2+b^2}$ 的大小。

\medskip
\textbf{二、为什么偏偏是 $a^2k^2+b^2$？}\par
把椭圆上的点写为 $P(a\cos\theta,b\sin\theta)$。对于斜率为 $k$ 的直线，关键不在 $x$ 或 $y$ 各自多大，而在
\[
y-kx=-ak\cos\theta+b\sin\theta.
\]
右边是两个正弦、余弦分量的组合，它能取得的\textbf{最大绝对值}就是
\[
\sqrt{(-ak)^2+b^2}=\sqrt{a^2k^2+b^2}.
\]
直线 $y=kx+m$ 要经过椭圆上的某个点，就要求 $m=y-kx$。当 $m$ 刚好达到这个最大值或最小值，直线便由“穿过”转为“刚好碰到”，于是得到相切公式。

\medskip
\textbf{三、几何图像：把椭圆拉伸回圆}\par
令 $X=x/a,\ Y=y/b$，椭圆就变成单位圆 $X^2+Y^2=1$，而原直线变成
\[
-akX+bY=m.
\]
在新坐标平面中，\textbf{圆心到直线的距离}是
\[
d=\frac{|m|}{\sqrt{a^2k^2+b^2}}.
\]
圆的半径是 $1$，故 $d<1$ 表示相交，$d=1$ 表示相切，$d>1$ 表示相离。\textbf{记住：先伸缩为圆，再用点到直线的距离；不要把新坐标中的距离误当作原椭圆平面里的距离。}

\medskip
\textbf{四、三个最快识别信号}\par
\begin{enumerate}
    \item \textbf{出现“斜率已知、求切线”}：直接写 $y=kx\pm\sqrt{a^2k^2+b^2}$。
    \item \textbf{出现“直线与椭圆相切”}：先写 $m^2=a^2k^2+b^2$，一般不必联立方程算判别式。
    \item \textbf{出现“过定点引切线”}：设 $y-y_0=k(x-x_0)$，即 $m=y_0-kx_0$，代入
    \[
    (y_0-kx_0)^2=a^2k^2+b^2,
    \]
    解 $k$，最后\textbf{另查竖直切线} $x=\pm a$ 是否经过该点。
\end{enumerate}

\textbf{一句话理解}：$\sqrt{a^2k^2+b^2}$ 不是随意拼出的根号，而是\textbf{椭圆上 $y-kx$ 的最大绝对值}；达到极限时，平行移动的直线刚好与椭圆相切。
\end{explanationbox}
